\chapter{Eight Leggy Essay}

\section{The History of the Chinese Eight-Legged Essay}

The eight-legged essay was a rigid writing format used in the Chinese imperial examination system from the Ming dynasty onward. It was designed to test candidates on their knowledge of the Confucian classics and their ability to follow a prescribed structure. The name comes from the essay's eight sections, each with a specific function and style, which candidates had to follow precisely. Mastery of this form became essential for anyone seeking a government position.

Over time, the eight-legged essay evolved from a tool for testing knowledge into an art of formula. Students memorized patterns, stock phrases, and model essays. Success depended less on original thought and more on reproducing the correct structure and approved interpretations of the classics. A talented candidate was not necessarily the one who understood the world best, but the one who could fit words into the required mold most smoothly.

By the late imperial period, criticism of the eight-legged essay was widespread. Reformers argued that it stifled independent thinking and trapped the educated elite in a closed world of textual games. Yet the system persisted because it was predictable, controllable, and familiar. Generations of scholars devoted their lives to perfecting a writing style that rewarded conformity over creativity.

When the Opium Wars and subsequent foreign invasions shook China in the nineteenth century, the limitations of this system became painfully clear. The country faced new technologies, new weapons, and new political ideas, but its educated class had been trained to excel at writing essays within a fixed pattern, not at thinking for themselves about unfamiliar problems. When real change arrived, almost no one in the official system knew how to respond. The eight-legged essay had produced experts in form, not in judgement.

\section{Eight-Legged Essays in Annual Reports}

There are actually quite a lot of eight-legged essays in annual reports. Companies fill the contents like the ancient Chinese students did: they follow prescribed formats, use stock phrases, and reproduce approved structures. Annual reports have become formulaic documents where companies slot in the required sections—executive summary, financial highlights, risk factors, corporate governance—without necessarily engaging in independent thinking about their actual business situation.

However, the business world is brutal and requires independent thinking. Markets change, technologies disrupt, competitors emerge, and customer preferences shift. Companies that simply fill in the template, that reproduce the same language year after year, that follow the formula without understanding what it means, are unlikely to survive. Like the Qing dynasty, which collapsed when it could no longer adapt to a changing world, companies that write eight-legged essays instead of thinking independently about their challenges will eventually fail. The form may look correct, but the substance is missing, and in business, substance is what matters.

\section{Business World Eight Leggy Essays}
We will present some excerpts of these eight-legged essays in annual reports.

\begin{quote}
    Together with our customers, we not only rise to meet today’s challenges, but we are also shaping the future by empowering each other to continuously improve and deliver better outcomes. Technology is essential in addressing the most profound economic, environmental, and social challenges. Software and analytics help predict and prevent disasters, make cities smarter, and foster equity in workplaces.
\end{quote}
    
It talks about everything, yet still saying nothing. What is the value that the software brings to the customers?

\begin{quote}
    The Group’s aspirations are subject to various external and internal factors, some of which it cannot influence. Successful
achievement of the bank’s 2025 strategic targets may be adversely impacted by reduced revenue generating capacities
of some of the bank’s core businesses should downside risks crystallize. These risks include but are not limited to the
challenging macroeconomic environment in Europe, in particular Germany, and the potential for a re-emergence of
inflationary pressures impacting the interest rate outlook. A number of geopolitical trends may exacerbate these risks,
including potential for widespread imposition of trade tariffs by the new U.S. administration.
\end{quote}

The challenge applies to all the international companies and you forget to mention climate change and Ukraine war.
And come on, you don't need so many words to say it. What is the unique risk for the company?

\begin{quote}
    What a remarkable year it has been! Put simply, 2024 was outstanding for Deutsche Telekom: all of our key financial performance indicators clearly exceeded market expectations. In November, for the first time in over two decades, our share surpassed the 30 euro mark.
\end{quote}

Focus on scoreboard is bad, focus on share price is worse. I understand your bonus hangs on the share 
price and you don't need to tell me. By the way, around 2000 the share price 
of Telekom has surpassed 100 euro, afterwards there was a crash. So what is it to be proud of?

\begin{quote}
    We have a strong foundation for future success. Building on our position of 
strength as a leading technology company, Siemens has launched the 
ONE Tech Company program to achieve the next level of performance and value 
creation. The program aims to ensure that the company leverages the 
opportunities arising from the historic market shifts and from technological 
disruptions. 
\end{quote}

Lots of self congratulations: strong, strength, next level. Removing these terms, nothing is left.


